\chapter{Diagnosing Multistability Versus Chaos}
\label{ch:diagnosing-multistability-chaos}

The same raster plot can be explained by two very different mechanisms. In a
multistable network, the deterministic large-network dynamics contains stable
attractors, and finite-size fluctuations cause rare escapes between their
basins. In a chaotic mesoscale network, the large-network dynamics itself has no
stable attractor on the relevant timescale; nearby cluster-rate trajectories
separate because the deterministic cluster-level vector field is unstable.

This chapter develops diagnostics that distinguish these mechanisms. The goal
is not to label every irregular trace as ``chaotic''. The goal is to ask which
source of variability is necessary for the switching to persist.

\section*{Learning Goals}

After this chapter you should be able to:
\begin{itemize}
    \item use the \(\xi\)-\(\kappa\) plane to separate finite-size noise from
    quenched heterogeneity;
    \item predict how switch times should scale in finite-size multistability
    versus mesoscopic chaos;
    \item interpret perturbation recovery, Lyapunov exponents, and zero-lag
    autocorrelation shape as dynamical diagnostics;
    \item explain why fully connected regime C is the clean many-cluster limit
    for SCS-like macroscopic chaos without finite-size effects.
\end{itemize}

\section{From Spikes to Diagnostic Time Series}
\label{sec:ch09-spikes-to-diagnostics}

The diagnostics in this chapter should be computed from a common data product,
not separately redefined for each figure. For each simulation store the
connectivity seed, stochastic seed convention, cluster labels, cell types,
spike times, filtered population rates, and, when available, the four input
current components defined in \cref{sec:ch09-transition-triggered}. The
metadata must record local counts \(n_E,n_I\) and total counts
\(N_E^{\mathrm{tot}}=Qn_E\), \(N_I^{\mathrm{tot}}=Qn_I\) separately. In this
chapter \(\xi=1/\sqrt{n_E}\) when the size sweep changes the number \(n_E\) of
excitatory neurons inside one cluster.

\begin{keyequation}[Spike binning and cluster-rate filtering]
Let \(t_m=m\Delta_b\), and let \(n_\alpha\) be the number of neurons of type
\(\alpha\) inside one cluster. The binned population activity and a causal
exponential rate estimate are
\begin{equation}
    B_{a,m}^{\alpha}
    =
    \frac{1}{n_\alpha\Delta_b}
    \sum_{i=1}^{n_\alpha}
    \int_{t_m}^{t_m+\Delta_b}
    s_{a,i}^{\alpha}(t)\,dt ,
    \qquad
    r_{a,m}^{\alpha}
    =
    \lambda_f r_{a,m-1}^{\alpha}
    +
    (1-\lambda_f)B_{a,m}^{\alpha},
    \label{eq:ch09-discrete-rate-filter}
\end{equation}
with \(\lambda_f=\exp(-\Delta_b/\tau_f)\). The units of \(B\) and \(r\) are
spikes per neuron per unit time. A noncausal Gaussian smoother is acceptable
for offline state detection, but the filter width must be reported because it
sets the shortest resolvable transition time.
\end{keyequation}

\paragraph{State detection protocol.}
Cluster states are binary summaries of the filtered rates, not direct spike
counts. A reproducible detector has the following steps.
\begin{enumerate}
    \item Discard burn-in and compute \(r_a^E(t)\) on a fixed bin grid for all
    clusters. Use the same filter for every cluster and every network size in a
    scaling comparison.
    \item Choose an activity threshold \(r_0\) from the valley between inactive
    and active modes of the pooled cluster-rate distribution. If the
    distribution is not clearly bimodal, use a prespecified quantile or
    mixture-model rule and report the sensitivity to the choice.
    \item Use hysteresis thresholds \(r_{\mathrm{on}}>r_{\mathrm{off}}\) around
    \(r_0\). A cluster turns on only after crossing \(r_{\mathrm{on}}\) and
    turns off only after crossing \(r_{\mathrm{off}}\). Merge events separated
    by less than a minimum dwell time \(T_{\min}\), typically several rate
    filter constants.
    \item Define the active-state vector
    \[
        S_a(t)=\mathbf 1\{r_a^E(t)\text{ is in the active hysteresis state}\},
        \qquad
        M(t)=\sum_{a=1}^Q S_a(t).
    \]
    Lifetimes are maximal intervals on which \(S_a(t)=1\). Transitions are
    threshold crossings that survive the dwell-time rule.
    \item For transition-triggered averages, keep an event only when no other
    cluster changes state in the exclusion window. The exclusion width should
    be at least the larger of the filter width and the median transition
    duration.
\end{enumerate}

\section{Two Mechanisms With Similar Traces}
\label{sec:ch09-two-mechanisms}

Let \(r(t)=(r_1^E,r_1^I,\ldots,r_Q^E,r_Q^I)\) denote the vector of cluster
rates. A schematic mesoscale model is
\begin{equation}
    \frac{d r}{dt}
    =
    F(r;\kappa,G)
    +
    \xi\,\eta(t),
    \qquad
    \xi=\frac{1}{\sqrt{n_E}}.
    \label{eq:ch09-schematic-mesoscale}
\end{equation}
The fixed random object \(G\) is the quenched inter-cluster disorder; its
amplitude is \(\kappa\). The process \(\eta(t)\) represents finite-size
population noise. At fixed local E/I fractions, let
\(n_c=n_E+n_I\), \(n_E=\gamma n_c\), and \(n_I=(1-\gamma)n_c\), so
\(\xi=(\gamma n_c)^{-1/2}\propto n_c^{-1/2}\). Throughout this chapter,
\(\kappa\) is the scalar heterogeneity axis; bare \(g\) is reserved for
inhibitory balance gains, and channel-specific quenched weight spreads are
written as \(\sigma_{\alpha\beta}^{\mathrm{het}}\). The same model therefore
contains two conceptually distinct limits:
\begin{equation}
    \text{finite-size multistability: }
    F \text{ has stable attractors and } \xi>0,
    \label{eq:ch09-multistable-limit}
\end{equation}
and
\begin{equation}
    \text{mesoscopic chaos: }
    F \text{ has a chaotic attractor and } \xi=0,\ \kappa>0 .
    \label{eq:ch09-chaotic-limit}
\end{equation}
In the first case, switching is a stochastic escape phenomenon. In the second,
switching is generated by the deterministic dynamics of the cluster rates.

\section[The xi-kappa Plane]{The \(\xi\)-\(\kappa\) Plane}
\label{sec:ch09-xi-kappa-plane}

The cleanest organizing picture is the plane with finite-size amplitude
\(\xi\) on one axis and quenched heterogeneity \(\kappa\) on the other.
Increasing the number of excitatory neurons per cluster, equivalently
increasing \(n_c\) at fixed \(\gamma\), moves the network toward \(\xi=0\).
Increasing the variability of inter-cluster weights or inter-cluster
connectivity moves the network upward in \(\kappa\).

\begin{table}[t]
\centering
\small
\begin{tabularx}{\textwidth}{@{}p{0.18\textwidth}p{0.35\textwidth}X@{}}
\toprule
\textbf{Region} & \textbf{Operational move} & \textbf{Interpretation} \\
\midrule
Large \(\xi\), small \(\kappa\) & Use small \(n_E\) and nearly homogeneous inter-cluster couplings & Switching can be dominated by finite-size escapes between stable cluster states. \\
Small \(\xi\), small \(\kappa\) & Increase \(n_E\) while keeping heterogeneity weak & Stable attractors should stop switching if finite-size noise was the only drive. \\
Large \(\xi\), large \(\kappa\) & Combine small clusters with frozen disorder & Irregular traces may mix finite-size hopping and deterministic heterogeneity effects. \\
Small \(\xi\), large \(\kappa\) & Increase \(n_E\) while retaining quenched inter-cluster disorder & Persistent switching here is the target evidence for mesoscopic chaos. \\
\bottomrule
\end{tabularx}
\caption{The \(\xi\)-\(\kappa\) diagnostic plane. The horizontal move is a
local-size sweep, \(\xi=1/\sqrt{n_E}\); the vertical move changes quenched
inter-cluster heterogeneity without changing the mean input.}
\label{tab:ch09-xi-kappa-plane}
\end{table}

This plane is useful because it makes qualitative claims falsifiable. A trace
that looks metastable at one network size is not enough. If the mechanism is
finite-size escape, the switching time should grow as \(\xi\) decreases. If the
mechanism is mesoscopic chaos, the switching time should approach a nonzero
large-network value controlled by \(\kappa\).

\section{Size Scaling and Switch-Time Plateaus}
\label{sec:ch09-switch-time-plateaus}

Let \(T_{\mathrm{sw}}\) be a characteristic switching time, such as the mean
cluster lifetime or the inverse transition rate between active-cluster
configurations. In a metastable attractor picture, a switch requires a rare
fluctuation across a basin boundary. A Kramers-like scaling gives the intuition
\begin{equation}
    T_{\mathrm{sw}}(\xi,\kappa\approx 0)
    \sim
    \exp\!\left(\frac{\Delta}{\xi^2}\right),
    \label{eq:ch09-kramers-scaling}
\end{equation}
where \(\Delta\) is an effective barrier. The equation is not meant to assert
that every clustered spiking network has a literal scalar potential. It states
the key scaling: if the escape is driven by finite-size noise, reducing
\(\xi\) should rapidly lengthen the switch time.

In a chaotic mesoscale regime, the large-network vector field itself produces
irregular motion. Then
\begin{equation}
    \lim_{\xi\to 0}T_{\mathrm{sw}}(\xi,\kappa)
    =
    T_{\mathrm{chaos}}(\kappa)
    <\infty .
    \label{eq:ch09-plateau}
\end{equation}
The practical test is therefore a plateau test. Increase \(n_c\), hence
\(n_E=\gamma n_c\), while holding the cluster-level parameters and quenched
realization statistics fixed. If the switching time first increases but then
saturates, finite-size noise may have been masking an underlying chaotic
process. If switching simply disappears, the evidence points to multistable
attractors destabilized by finite-size fluctuations.

\section{Perturbation Recovery Versus Divergence}
\label{sec:ch09-perturbation-diagnostics}

Perturbation experiments ask whether a small displacement is forgotten. Take
two copies of the same network realization, start them from the same
cluster-rate state, and perturb one copy by a small vector \(\delta r(0)\). A
natural distance is
\begin{equation}
    d(t)
    =
    \left[
    \frac{1}{Q}
    \sum_{a=1}^{Q}
    \left(
    \left(r_{a,1}^E(t)-r_{a,2}^E(t)\right)^2
    +
    \left(r_{a,1}^I(t)-r_{a,2}^I(t)\right)^2
    \right)
    \right]^{1/2}.
    \label{eq:ch09-trajectory-distance}
\end{equation}
The subscripts \(1\) and \(2\) label the two trials. In a metastable basin, a
small perturbation should decay on a local relaxation timescale:
\begin{equation}
    d(t) \approx d(0)e^{-t/\tau_{\mathrm{relax}}},
    \label{eq:ch09-perturbation-decay}
\end{equation}
with \(\tau_{\mathrm{relax}}\) comparable to the single-neuron or membrane
timescale, not the much longer cluster lifetime. Only perturbations large
enough to cross a basin boundary, typically comparable to the finite-size
escape fluctuation, should produce a different active-cluster state. Thus the
metastable signature is a fast return for small perturbations and a thresholded
state change for large perturbations.

In a chaotic mesoscale regime, an arbitrarily small generic perturbation
initially grows:
\begin{equation}
    d(t) \approx d(0)e^{\lambda_{\max}t},
    \qquad \lambda_{\max}>0.
    \label{eq:ch09-perturbation-growth}
\end{equation}
The growth eventually saturates because rates are bounded, but the early
exponential window is the diagnostic. The timescale discriminator is as
important as the sign of the response: chaotic responses persist or diverge on
the cluster autocorrelation timescale \(\tau_c\), whereas metastable recovery
from a subthreshold perturbation is fast, on the order of \(\tau_m\) or the
finite-size relaxation time.

\begin{table}[t]
\centering
\small
\begin{tabularx}{\textwidth}{@{}p{0.30\textwidth}p{0.30\textwidth}X@{}}
\toprule
\textbf{Perturbation size} & \textbf{Stable-basin response} & \textbf{Chaotic mesoscale response} \\
\midrule
Infinitesimal rate perturbation & Decays with local relaxation time & Grows during a linear exponential window. \\
Finite but subthreshold spiking perturbation & Recovers to the same threshold-defined state under common noise & Separates reproducibly before saturation if projected instability is present. \\
Basin-crossing perturbation & Produces a thresholded jump to another attractor & Not the diagnostic event; chaos is tested before saturation and before state labels diverge. \\
\bottomrule
\end{tabularx}
\caption{Perturbation signatures. Basin crossing is not enough to claim chaos;
the key test is reproducible growth of small perturbations under deterministic
or common-noise conditions.}
\label{tab:ch09-perturbation-signatures}
\end{table}

In spiking simulations one must be careful about stochasticity. The cleanest
version compares either deterministic mesoscale trajectories or spiking trials
with controlled randomness so that the measured separation is caused by the
perturbation rather than by unrelated random seeds. The quenched connectivity
realization should be identical in the two copies.

\section{Lyapunov Diagnostics}
\label{sec:ch09-lyapunov-diagnostics}

The maximal Lyapunov exponent is the asymptotic rate at which an infinitesimal
perturbation grows:
\begin{equation}
    \lambda_{\max}
    =
    \lim_{t\to\infty}
    \frac{1}{t}
    \log
    \frac{\|\delta r(t)\|}{\|\delta r(0)\|}.
    \label{eq:ch09-lyapunov}
\end{equation}
For a smooth mesoscale system \(dr/dt=F(r)\), the perturbation evolves by the
tangent equation
\begin{equation}
    \frac{d}{dt}\delta r(t)
    =
    D F(r(t))\,\delta r(t),
    \label{eq:ch09-tangent-equation}
\end{equation}
where \(D F\) is the Jacobian of the cluster-rate vector field along the
trajectory. A positive \(\lambda_{\max}\) is not the same as irregular-looking
activity; it is a mathematical statement about local exponential separation.

\paragraph{Deterministic mesoscopic Benettin algorithm.}
For an effective deterministic cluster model, write the full numerical state as
\(x(t)\). This state may include excitatory and inhibitory rates, adaptation
variables, synaptic filters, refractory or age-density bins, and any auxiliary
variables needed to make the dynamics Markovian. For the largest exponent,
integrate
\[
    \dot x=F(x),
    \qquad
    \dot v=D F(x)v
\]
with \(\|v(0)\|=1\). Every \(\Delta_R\) units of time, record
\(g_m=\|v(m\Delta_R)\|\), renormalize \(v\leftarrow v/g_m\), and continue.
After discarding an initial transient,
\begin{equation}
    \widehat\lambda_{\max}(T)
    =
    \frac{1}{M\Delta_R}
    \sum_{m=1}^{M}\log g_m .
    \label{eq:ch09-benettin-max}
\end{equation}
For a spectrum, evolve a matrix \(U\) of tangent vectors and perform a QR
factorization \(U=QR\) at each renormalization time:
\begin{equation}
    \widehat\lambda_j(T)
    =
    \frac{1}{M\Delta_R}
    \sum_{m=1}^{M}\log |R_{jj}^{(m)}| .
    \label{eq:ch09-qr-lyapunov}
\end{equation}
The renormalization interval \(\Delta_R\) should be long enough that the
integrator is not just measuring one time step and short enough that the
perturbation has not saturated; in practice it is tested over a small grid and
kept well below the broad cluster autocorrelation time.

\paragraph{Finite-difference shadowing when the Jacobian is unavailable.}
Run two deterministic mesoscopic trajectories with identical parameters and
initial separation \(\varepsilon v\). Evolve for \(\Delta_R\), estimate the
growth factor
\[
    g_m=\frac{\|x_2(m\Delta_R)-x_1(m\Delta_R)\|}{\varepsilon},
\]
then reset the second trajectory to
\[
    x_2\leftarrow x_1
    +
    \varepsilon
    \frac{x_2-x_1}{\|x_2-x_1\|}.
\]
Accumulating \(\log g_m\) gives the same maximal-exponent estimate as
\eqref{eq:ch09-benettin-max} when the perturbation remains in the linear
regime. Repeat the estimate for at least two values of \(\varepsilon\); a
linear-regime estimate should not change when \(\varepsilon\) is made smaller.

\paragraph{Common-noise spiking perturbation protocol.}
For spiking simulations, independent random seeds measure noise divergence, not
deterministic sensitivity. A controlled protocol is:
\begin{enumerate}
    \item Generate one quenched connectivity realization and one complete
    external-noise or stochastic-threshold stream that can be replayed.
    \item Clone the network state after burn-in, including membrane potentials,
    adaptation variables, refractory states, and synaptic filter states.
    \item Perturb the clone by one member of a calibrated ladder: one membrane
    potential by \(\varepsilon\), all membrane potentials in one cluster by
    \(\varepsilon\), one added spike, or one added spike in every neuron of one
    cluster.
    \item Replay the same noise streams in both copies and measure the
    filtered-rate distance \(d(t)\) from \eqref{eq:ch09-trajectory-distance}.
    \item Estimate slopes of \(\log d(t)\) only before saturation and before
    the two trajectories land in clearly different threshold-defined states.
\end{enumerate}
This gives a random-dynamical-system or projected mesoscale exponent. It is
strong evidence only when the estimate is reproducible under the perturbation
ladder and remains positive as the finite-size amplitude is reduced.

\paragraph{Stochastic caveats and convergence checks.}
Finite-size Langevin equations with multiplicative noise can have positive
common-noise exponents even when the deterministic drift is stable. Therefore
the claim ``mesoscopic chaos'' should be reserved for a positive exponent in
the deterministic large-cluster theory or for a positive common-noise exponent
that approaches a nonzero limit as \(\xi\to 0\). Report convergence plots of
\(\widehat\lambda_{\max}(T)\), vary the time step, renormalization interval,
norm, transient length, and perturbation size, and verify the implementation on
a known stable linear system and on a reference chaotic rate model before using
it as evidence in the clustered spiking model.

\section{Autocorrelation Diagnostics and the Zero-Lag Kink}
\label{sec:ch09-autocorrelation-kink}

Let
\begin{equation}
    C_c(\tau)
    =
    \left\langle
    \delta r_c(t)\,\delta r_c(t+\tau)
    \right\rangle_t,
    \qquad
    \delta r_c(t)=r_c(t)-\langle r_c\rangle_t ,
    \label{eq:ch09-cluster-autocorrelation}
\end{equation}
be the autocorrelation of a cluster-rate trace. Finite-size fluctuations often
contribute a fast Ornstein-Uhlenbeck-like term,
\begin{equation}
    C_{\mathrm{fs}}(\tau)
    =
    A_\xi e^{-|\tau|/\tau_\xi}.
    \label{eq:ch09-ou-autocorrelation}
\end{equation}
This function has a cusp at zero lag:
\begin{equation}
    \left.\frac{d C_{\mathrm{fs}}}{d\tau}\right|_{\tau=0^+}
    =
    -\frac{A_\xi}{\tau_\xi},
    \qquad
    \left.\frac{d C_{\mathrm{fs}}}{d\tau}\right|_{\tau=0^-}
    =
    \frac{A_\xi}{\tau_\xi}.
    \label{eq:ch09-kink-derivative}
\end{equation}
By contrast, the autocorrelation of a differentiable deterministic chaotic rate
process is even and smooth at zero lag, so its first derivative vanishes there.
An observed autocorrelation can be decomposed schematically as
\begin{equation}
    C(\tau)
    =
    C_{\mathrm{chaos}}(\tau)
    +
    A_\xi e^{-|\tau|/\tau_\xi}.
    \label{eq:ch09-autocorrelation-mixture}
\end{equation}
The kink amplitude \(A_\xi/\tau_\xi\) estimates the fast finite-size
contribution. It should shrink with \(n_E\). A smooth zero-lag autocorrelation
that retains a broad timescale as \(n_E\to\infty\) is consistent with
mesoscopic chaos.

The diagnostic contrast can be read directly from the local expansion:
\[
    \rho_{\mathrm{fs}}(\tau)=1-k_\xi|\tau|+O(\tau^2),
    \qquad
    \rho_{\mathrm{smooth}}(\tau)=1-\frac{1}{2}b\tau^2+O(\tau^4).
\]
The first form has a nonzero one-sided slope at the origin; the second has a
vanishing first derivative. A measured trace may contain both pieces, which is
why the fitted \(k_\xi\) must be tracked over the \(n_E\) size ladder rather
than interpreted from one network size.

This diagnostic is mathematically clean but experimentally delicate. It requires
enough temporal resolution to see the zero-lag shape and enough trials or
stationary data to estimate the autocorrelation reliably.

\paragraph{Operational estimator.}
For a discrete filtered rate \(r_m=r(m\Delta_b)\), remove burn-in and estimate
the autocovariance with valid-overlap normalization:
\begin{equation}
    \widehat C[\ell]
    =
    \frac{1}{M-\ell}
    \sum_{m=0}^{M-\ell-1}
    (r_m-\overline r)(r_{m+\ell}-\overline r),
    \qquad
    0\le \ell<M .
    \label{eq:ch09-discrete-autocovariance}
\end{equation}
FFT implementations are fine, but they must divide by \(M-\ell\), not by a
constant \(M\), because large lags have fewer overlapping samples. Normalize by
\(\widehat\rho[\ell]=\widehat C[\ell]/\widehat C[0]\). The HWHM is obtained by
linear interpolation between the first adjacent lags satisfying
\(\widehat\rho[\ell]>1/2\) and \(\widehat\rho[\ell+1]\le 1/2\):
\begin{equation}
    \widehat\tau_{\mathrm{HWHM}}
    =
    \Delta_b
    \left[
        \ell
        +
        \frac{\widehat\rho[\ell]-1/2}
        {\widehat\rho[\ell]-\widehat\rho[\ell+1]}
    \right].
    \label{eq:ch09-hwhm-interpolation}
\end{equation}
If the autocorrelation does not cross one half inside the recording window,
report a right-censored lower bound instead of extrapolating.

For the zero-lag kink, discard the raw same-bin spike-count peak and use
filtered rates. Fit the symmetric small-lag model
\begin{equation}
    \rho(\tau)
    \approx
    1-k_{\xi}|\tau|
    +
    \frac{1}{2}b\tau^2,
    \qquad 0<|\tau|\le w ,
    \label{eq:ch09-kink-fit}
\end{equation}
over several windows \(w\). The finite-size prediction is that \(k_\xi\)
decreases with local population size, while a smooth deterministic rate process
has \(k_\xi\to 0\) as the bin width and fit window are refined. Block bootstrap
or network-realization bootstrap intervals should be shown because short traces
can manufacture apparent zero-lag slopes.

\section{Transition-Triggered Diagnostics}
\label{sec:ch09-transition-triggered}

Suppose a cluster \(c\) turns on at times \(t_m^\uparrow\) and turns off at
times \(t_m^\downarrow\), detected by upward and downward crossings of a rate
threshold \(r_0\). Use isolated transitions: no other cluster should cross
threshold inside a window
\([t_m-T_C,t_m+T_C]\), where \(T_C\) is a typical cluster lifetime. This
selection prevents a triggered average for cluster \(c\) from being dominated
by a nearly simultaneous transition elsewhere.

For an observable \(X(t)\), define the transition-triggered average
\begin{equation}
    \operatorname{TTA}_X(\tau)
    =
    \frac{1}{M}
    \sum_{m=1}^{M}
    X(t_m^\uparrow+\tau).
    \label{eq:ch09-tta}
\end{equation}
The same definition is used for \(t_m^\downarrow\) when diagnosing on-to-off
events.

The input current to an excitatory neuron in the focus cluster should be split
by source location and source type,
\begin{equation}
    I_{c,i}^{E}(t)
    =
    I_{\mathrm{in},i}^{E}(t)
    +
    I_{\mathrm{in},i}^{I}(t)
    +
    I_{\mathrm{out},i}^{E}(t)
    +
    I_{\mathrm{out},i}^{I}(t).
    \label{eq:ch09-input-split}
\end{equation}
Here ``in'' means presynaptic neurons inside cluster \(c\), and ``out'' means
presynaptic neurons in other clusters. For \(x\in\{\mathrm{in},\mathrm{out}\}\)
and \(\beta\in\{E,I\}\), define the spatial input-current variance across the
excitatory neurons of the focus cluster by
\begin{equation}
    \left(\sigma_x^\beta(t)\right)^2
    =
    \frac{1}{n_E}
    \sum_{i=1}^{n_E}
    \left(
        I_{x,i}^{\beta}(t)
        -
        \overline I_x^\beta(t)
    \right)^2 .
    \label{eq:ch09-sigma-components}
\end{equation}
The four components
\(\sigma_{\mathrm{in}}^E,\sigma_{\mathrm{in}}^I,
\sigma_{\mathrm{out}}^E,\sigma_{\mathrm{out}}^I\) ask whether transitions are
preceded by fluctuations from local excitation, local inhibition,
inter-cluster excitation, or inter-cluster inhibition. The corresponding mean
currents \(\overline I_x^\beta(t)\) should be triggered in the same way. If
finite-size fluctuations drive transitions, one expects one or more
\(\sigma_x^\beta\) components to rise immediately before \(t_m\). If inhibition
controls the escape, the pre-transition signature may appear as a decrease in
inhibitory mean current rather than as an increase in excitation.

The roadmap also calls for a directional state-transition average. Condition on
a specific transition \(A\to B\), where \(r_A\) and \(r_B\) are the cluster-rate
vectors for the two macroscopic states, and define
\begin{equation}
    e_{AB}
    =
    \frac{r_B-r_A}{\|r_B-r_A\|}.
    \label{eq:ch09-transition-direction}
\end{equation}
The rate displacement decomposes into a transition coordinate and an orthogonal
coordinate,
\begin{equation}
    z_{\parallel}(t)
    =
    e_{AB}^{\mathsf T}\bigl(r(t)-r_A\bigr),
    \qquad
    z_{\perp}^2(t)
    =
    \left\|r(t)-r_A\right\|^2-z_{\parallel}^2(t).
    \label{eq:ch09-transition-projections}
\end{equation}
A scalar variance burst averaged over all off-to-on events is not enough. For a
specific \(A\to B\) transition, the diagnostic asks whether the triggered mean
rate moves along \(r_B-r_A\) before the crossing, while fluctuations orthogonal
to that direction stay small or decrease.

The complementary cross-correlation asks whether fluctuations in one quantity
lead changes in another:
\begin{equation}
    C_{r,\sigma^2}(\tau)
    =
    \left\langle
    \bigl(r_c(t)-\overline r_c\bigr)
    \bigl(\sigma_c^2(t+\tau)-\overline{\sigma_c^2}\bigr)
    \right\rangle_t .
    \label{eq:ch09-cross-correlation}
\end{equation}
Here \(\sigma_c^2(t)\) can be replaced by any of
\((\sigma_{\mathrm{in}}^E)^2,(\sigma_{\mathrm{in}}^I)^2,
(\sigma_{\mathrm{out}}^E)^2,(\sigma_{\mathrm{out}}^I)^2\). A peak at negative
lag means that input-current variance rises before the rate transition; a peak
at positive lag means variance follows the rate change. These observables do
not by themselves prove chaos, but they identify the mechanism that triggers
state changes.

\paragraph{Transition-triggered Fano factors.}
For a focus population \(\mathcal P\), event time \(t_m\), offset \(\tau\), and
counting window \(W\), define
\begin{equation}
    N_m^{\mathcal P}(\tau;W)
    =
    \sum_{i\in\mathcal P}
    \int_{t_m+\tau}^{t_m+\tau+W}s_i(t)\,dt .
    \label{eq:ch09-transition-counts}
\end{equation}
The transition-triggered Fano curve is
\begin{equation}
    F_{\mathcal P}^{\mathrm{TTA}}(\tau;W)
    =
    \frac{\operatorname{Var}_m[N_m^{\mathcal P}(\tau;W)]}
    {\mathbb E_m[N_m^{\mathcal P}(\tau;W)]}.
    \label{eq:ch09-transition-fano}
\end{equation}
Compute it separately for off-to-on and on-to-off events, and compare it with
rate-matched control windows that have similar mean rate but are not near a
state transition. Without this control, a Fano increase can be a trivial
consequence of a rising mean rate rather than a pre-transition fluctuation
signature.

\paragraph{Transition-triggered dimensionality.}
Let \(r(t)\in\mathbb R^Q\) be the vector of excitatory cluster rates, or
\(\mathbb R^{2Q}\) if E and I rates are concatenated. Around transitions,
there are two useful covariance matrices:
\begin{equation}
    C_{\mathrm{raw}}(\tau)
    =
    \operatorname{Cov}_m[r(t_m+\tau)],
    \qquad
    C_{\mathrm{res}}(\tau)
    =
    \operatorname{Cov}_m[
    r(t_m+\tau)-\operatorname{TTA}_r(\tau)].
    \label{eq:ch09-transition-covariances}
\end{equation}
The raw covariance reports the dimension of the whole transition cloud; the
residual covariance reports variability around the mean transition path. Use
\[
    d_{\mathrm{PR}}(\tau)
    =
    \frac{\operatorname{Tr}(C(\tau))^2}{\operatorname{Tr}(C(\tau)^2)}
\]
for either covariance. In a stereotyped attractor-to-attractor transition,
\(d_{\mathrm{PR}}\) should be small near the crossing after conditioning on
the transition \(A\to B\). A chaotic transition ensemble may retain a larger
residual dimension, especially away from the threshold crossing.

\section{Fully Connected Scaling Regimes A, B, and C}
\label{sec:ch09-abc-regimes}

The roadmap separates three large-network limits for fully connected
inter-cluster networks. Let \(Q\) be the number of clusters, \(n_c=n_E+n_I\)
the neurons per cluster, and \(N^{\mathrm{tot}}=Qn_c\) the total neuron count.
Inter-cluster weights have the template form
\begin{equation}
    J_{i_a^\alpha,j_b^\beta}
    =
    J_{\alpha\beta}^{\mathrm{base}}(Q,n_c)\,
    \check j_{ab}^{\alpha\beta},
    \qquad
    \check j_{ab}^{\alpha\beta}\sim
    \mathcal{N}\!\left(1,(\sigma_{\alpha\beta}^{\mathrm{het}})^2\right),
    \quad a\ne b.
    \label{eq:ch09-fully-connected-random-weights}
\end{equation}
All neurons in a given source-target cluster pair share the same pair factor
\(\check j_{ab}^{\alpha\beta}\). The base scale
\(J_{\alpha\beta}^{\mathrm{base}}(Q,n_c)\) is part of the model definition and
must be written in the manifest; it may differ between dense, fixed-indegree,
and mean-field templates.

The following expression is only a channel-ranking heuristic, not a universal
variance law. For one source cluster \(b\), the typical quenched spread of the
\(\beta\to\alpha\) contribution to a cluster-averaged input is proportional to
\begin{equation}
    H_{ab}^{\alpha\beta}
    =
    n_\beta
    \left(J_{\alpha\beta}^{\mathrm{base}}(Q,n_c)\right)^2
    (\sigma_{\alpha\beta}^{\mathrm{het}})^2
    r_b^\beta ,
    \label{eq:ch09-pair-heterogeneity-score}
\end{equation}
up to synaptic filtering, correlations, and the details of the connectivity
construction. The implementable diagnostic is empirical: store the four
outside-cluster current components, compute their within-realization temporal
variability and between-realization bias, and compare those quantities across
the same \(Q,n_E,n_I,\kappa\) manifest grid. Finite-size population noise is
tracked separately by \(\xi=1/\sqrt{n_E}\).

\paragraph{Regime A: many clusters with fixed cluster size.}
Here \(n_c\) is fixed and \(Q\to\infty\). The finite-size amplitude
\(\xi=1/\sqrt{n_E}\) does not vanish. Quenched variability is also present.
This regime is useful for studying switching, but it cannot by itself prove
that switching survives the removal of finite-size effects.

\paragraph{Regime B: few large clusters.}
Here \(Q\) is fixed and \(n_c\to\infty\). The finite-size amplitude vanishes, leaving
a deterministic low-dimensional cluster system with a finite number of random
couplings. The dynamics can be very sensitive to the specific draw of
\(\check j_{ab}^{\alpha\beta}\). This limit can show deterministic irregular
activity, but it does not yet realize a many-unit dynamic mean-field limit.

\paragraph{Regime C: many large clusters.}
Here both \(Q\to\infty\) and \(n_c\to\infty\), for example
\(Q=(N^{\mathrm{tot}})^\rho\) and
\(n_c=(N^{\mathrm{tot}})^{1-\rho}\) with \(0<\rho<1\). The finite-size
amplitude vanishes, but the number of random cluster couplings grows.
This is the regime that most directly realizes Sompolinsky-Crisanti-Sommers
style macroscopic chaos at the cluster level: many effective units, random
couplings, no finite-size sampling noise required.

\begin{table}[t]
\centering
\small
\begin{tabularx}{\textwidth}{@{}p{0.16\textwidth}p{0.28\textwidth}p{0.22\textwidth}X@{}}
\toprule
\textbf{Regime} & \textbf{Scaling} & \textbf{Finite-size noise} & \textbf{Use in diagnosis} \\
\midrule
A & \(Q\to\infty\), fixed \(n_c\) & Remains because \(n_E\) is fixed & Many-cluster statistics with finite-size hopping still present. \\
B & Fixed \(Q\), \(n_c\to\infty\) & Vanishes & Tests whether a particular finite cluster graph has deterministic irregular dynamics. \\
C & \(Q\to\infty\), \(n_c\to\infty\) & Vanishes & Clean target for many-unit cluster-level chaos and DMFT comparison. \\
\bottomrule
\end{tabularx}
\caption{A/B/C scaling regimes. The classification uses the manifest variables
\(Q,n_E,n_I\) and the empirical current decompositions, not an assumed closed
form for the input variance.}
\label{tab:ch09-abc-regimes}
\end{table}

\section{Size, Disorder, and Validation Protocols}
\label{sec:ch09-sweep-validation-protocols}

The diagnostic program is a set of controlled sweeps. The same analysis code
should run on spiking data, deterministic mesoscopic trajectories, and
synthetic validation traces.

\paragraph{Finite-size scaling protocol.}
Choose an architecture, a cluster-strength setting, an inhibitory-clustering
setting, and a quenched-disorder amplitude. Then run a size ladder
\(n_E\in\{n_1,n_2,\ldots\}\) inside each cluster while holding \(Q\), the
cluster-level weights, and the distribution of quenched couplings fixed. For
each \(n_E\), use several connectivity realizations and long recordings. The
primary response variable is a characteristic switch time:
\[
    T_{\mathrm{sw}}
    \in
    \{\text{mean active lifetime},\,
    \text{median active lifetime},\,
    \text{inverse transition rate}\}.
\]
Fit two deliberately simple alternatives:
\begin{align}
    \log T_{\mathrm{sw}}(n_E)
    &=
    a+b n_E
    \quad\text{finite-size escape template,}
    \notag\\
    T_{\mathrm{sw}}(n_E)
    &=
    T_\infty + \frac{b}{n_E}
    \quad\text{chaotic plateau template.}
    \label{eq:ch09-size-scaling-fits}
\end{align}
These are diagnostic templates, not exact laws. The finite-size interpretation
is credible only if the plateau fit fails and switches become right-censored
as \(n_E\) grows. The chaos interpretation is credible only if \(T_\infty\) is
positive, the same parameter point has a positive mesoscale Lyapunov estimate,
and the autocorrelation width is stable across the size ladder.

\paragraph{Quenched-disorder sweep protocol.}
At one or more large local population sizes, sweep the heterogeneity coordinate
\(\kappa\) while keeping mean weights, balance parameters, and depression rules
fixed. For each \(\kappa\), separate total rate variability into a
within-realization and a between-realization component:
\begin{equation}
    \operatorname{Var}_{t,G}(r_a)
    =
    \mathbb E_G[\operatorname{Var}_t(r_a\mid G)]
    +
    \operatorname{Var}_G(\mathbb E_t[r_a\mid G]),
    \label{eq:ch09-time-realization-variance-split}
\end{equation}
and, when input currents are stored, split the first term into
within-cluster E/I and outside-cluster E/I components. The first term measures
within-realization temporal dynamics; the second measures sample-to-sample
biases from the quenched draw. A genuine chaotic regime should not be reduced
to a few static high-rate clusters selected by the draw; it should show
within-realization temporal variability, positive instability, and broad state
exploration.

\paragraph{Network-generation pseudocode for diagnostic sweeps.}
The connectivity generator should be deterministic given a seed and should
return all random objects needed for reproducibility:
\begin{enumerate}
    \item assign neuron labels \((a,\alpha,i)\) with cluster \(a\), type
    \(\alpha\in\{E,I\}\), and within-cluster index \(i\);
    \item draw or construct fixed indegrees for every postsynaptic population
    block, separating within-cluster and outside-cluster sources;
    \item apply potentiation factors to within-cluster blocks and compensating
    depression factors to outside-cluster blocks;
    \item draw quenched pair factors or adjacency variables for outside-cluster
    blocks, enforcing the intended sign convention for E and I weights;
    \item write a connectivity manifest containing seed, indegrees, realized
    in-degree histograms, mean block weights, and block-weight variances.
\end{enumerate}
The manifest is part of the data, not bookkeeping. It is what allows the
diagnostic code to know whether a change in switching came from \(\xi\),
\(\kappa\), or an unintended drift in mean input.

\paragraph{Validation battery.}
Before analyzing research runs, the numerical pipeline should pass the
following tests.
\begin{enumerate}
    \item A homogeneous Poisson spike train has Fano factor near one after
    binning and a flat shift-corrected autocorrelation away from zero lag.
    \item A synthetic two-state telegraph process is recovered by the state
    detector with the correct lifetimes, transition times, and active-count
    distribution.
    \item An Ornstein-Uhlenbeck rate process produces the expected zero-lag
    cusp, and a smooth deterministic test signal produces \(k_\xi\approx 0\)
    under the kink fit.
    \item A known stable linear rate system gives \(\lambda_{\max}<0\), and a
    reference chaotic rate model gives a positive exponent with the Benettin or
    QR implementation.
    \item Replaying the same stochastic streams in two unperturbed spiking
    copies produces identical trajectories up to numerical precision; replaying
    independent streams is treated as a noise-control experiment, not a
    Lyapunov experiment.
    \item Varying the activity threshold, filter width, and minimum dwell time
    within a prespecified range does not change the qualitative phase label of
    a robust parameter point.
\end{enumerate}

\section{Finite-Size Effects Versus Mesoscopic Chaos}
\label{sec:ch09-decision-table}

The diagnostics are strongest when they agree:
\begin{itemize}
    \item If \(T_{\mathrm{sw}}\) diverges as \(\xi\to 0\), perturbations recover,
    the zero-lag autocorrelation has a shrinking finite-size kink, and
    \(\lambda_{\max}\le 0\), the evidence favors finite-size multistability.
    \item If \(T_{\mathrm{sw}}\) plateaus as \(\xi\to 0\), infinitesimal
    perturbations diverge, the broad autocorrelation persists without a
    finite-size kink, and \(\lambda_{\max}>0\), the evidence favors mesoscopic
    chaos.
\end{itemize}
Intermediate regimes are expected. A network may have stable basins at small
\(\kappa\), deterministic chaos at large \(\kappa\), and a mixed region where
finite-size noise and chaotic instability both influence transitions. The point
of the diagnostic program is to map that region rather than collapse it into a
single label.

\section{Chapter Summary}
\label{sec:ch09-summary}

The distinction between multistability and chaos is a distinction between
mechanisms, not appearances. The finite-size amplitude
\(\xi=1/\sqrt{n_E}=(\gamma n_c)^{-1/2}\) and quenched heterogeneity \(\kappa\)
organize the tests. A reproducible analysis starts by mapping spikes to
filtered cluster rates and threshold-defined state vectors with documented
filters, thresholds, and dwell-time rules. Network-size scaling asks whether
switch times plateau or diverge. Common-noise perturbation experiments and
deterministic Benettin or QR tangent algorithms ask whether nearby cluster-rate
trajectories recover or separate. Autocorrelation shape, especially the
zero-lag kink, separates fast finite-size noise from smooth deterministic rate
fluctuations. Transition-triggered Fano factors, dimensionality, current
variance decompositions, and cross-correlations identify which currents and
variances precede state changes. Regime C, with many large clusters, is the
clean mathematical limit for SCS-like macroscopic chaos without finite-size
effects, and the validation battery is what keeps this classification from
being an artifact of binning, thresholds, or random seeds.
